% titel: Blatt 5, A3: Radioaktiver Zerfall (C-14)
% skript: Kap. 2 (Exponentialfunktion), Kap. 3 (lim)

% Teil I: Werkzeuge, kein Lösungsweg.

Aus der Halbwertszeit die Zerfallskonstante $\lambda$ bestimmen und damit
das Alter einer Probe aus der verbliebenen Konzentration.

\begin{werkzeugkasten}[Warum ist der Zerfall exponentiell?]
\emph{Worum es geht:} Ein einzelner Atomkern \emph{vergisst} seine
Vorgeschichte: Zu jedem Zeitpunkt ist die Zerfallswahrscheinlichkeit pro
Zeiteinheit eine feste Konstante. Damit ist die \emph{Zerfallsrate}
proportional zur Zahl der noch vorhandenen Kerne $N$:
\begin{equation*}
  \frac{\dd N}{\dd t}=-\lambda N .
\end{equation*}
Das Minus zeigt (abnehmende Menge), $\lambda$ (die Zerfallskonstante) misst,
wie schnell. Die Differentialgleichung sagt: \emph{Änderung proportional zum
Bestand} -- und die einzige Funktion mit dieser Eigenschaft ist die
Exponentialfunktion. Über je ein gleiches Zeitintervall $\Delta t$ bleibt
derselbe Faktor $\mathrm e^{-\lambda\Delta t}$ übrig (gleiches
Wahrscheinlichkeitsgesetz). Deshalb:
\begin{equation*}
  N(t)=N_0\,\mathrm e^{-\lambda t}, \qquad \lambda>0 .
\end{equation*}
Die Halbwertszeit $T$ ist die Zeit, nach der die Hälfte übrig ist:
$N(T)=\tfrac12 N_0$, also $\mathrm e^{-\lambda T}=\tfrac12$, und damit
\begin{equation*}
  \lambda=\frac{\ln 2}{T}.
\end{equation*}
$\lambda$ hat die Einheit einer inversen Zeit (z. B. $\mathrm{1\,s^{-1}}$
bzw.\ eine inverse Jahres-Skala hier); $T$ eine Zeit.

\textit{Skript:} Kap. 2 (Exponentialfunktion).
\end{werkzeugkasten}

\begin{werkzeugkasten}[Vorgehen bei zwei Unbekannten]
Wir kennen $\lambda$ (aus $T$) und messen den Restanteil $q=N(t)/N_0$. Aus
$q=\mathrm e^{-\lambda t}$ folgt durch Logarithmieren (die \emph{Umkehrung}
der Exponentialfunktion \emph{löst} das $t$ aus dem Exponenten):
\begin{equation*}
  \ln q=-\lambda t
  \;\Rightarrow\;
  t=\frac{-\ln q}{\lambda}
  =\frac{\ln(1/q)}{\lambda}
  =\frac{T\,\ln(1/q)}{\ln 2}.
\end{equation*}
Für $q<1$ ist $\ln(1/q)>0$: das Alter kommt positiv heraus. Da sich $N_0$
im Quotienten $q$ weggekürzt hat, spielt die ursprüngliche Menge keine
Rolle -- es zählt nur der \emph{Anteil}, der noch da ist.

\textit{Skript:} Kap. 2, Kap. 3 (Stetigkeit von $\ln$).
\end{werkzeugkasten}

\begin{achtung}[Typische Fallen]
\begin{itemize}
  \item \textbf{Einheitensystem:} $\lambda$ und die Zeit müssen im selben
        Zeitsystem stehen (hier Jahre), sonst kommt ein unsinniges Alter
        heraus.
  \item \textbf{Prozentzahl statt Anteil:} $53\,\%$ ist $q=0{,}53$, nicht
        $q=53$. Das ist der Klassiker.
  \item \textbf{$\ln$ nicht vergessen:} Aus $q=\mathrm e^{-\lambda t}$
        wird $t$ durch $\ln$ isoliert; ohne $\ln$ bleibt die Zeit \emph{im
        Exponenten gefangen}.
  \item \textbf{$N_0$ verwechseln:} Es kürzt sich heraus; es zählt nur der
        gemessene Anteil $q$.
\end{itemize}
\end{achtung}
